Se aborda la detección y clasificación de especies en fotografías de cámaras
trampa del conjunto Caltech Camera Traps, donde el 70\% de las imágenes no
contiene ningún animal y el desbalance entre especies alcanza un factor 500. Se
ajusta un detector YOLO11 a 640 px sobre 12\,099 imágenes con bounding box; las
imágenes vacías no entran al entrenamiento ni a la evaluación, porque en
detección la ausencia de cajas ya expresa que no hay nada que detectar.
El modelo alcanza mAP@0.5 $= 0.7463$ (IC95\% $[0.7236, 0.7563]$) en ubicaciones
de cámara vistas durante el entrenamiento y $0.3978$ en ubicaciones nunca
vistas, una brecha de $0.3485$ con intervalos disjuntos. Controlando el cambio
de prevalencia de especies entre particiones, la brecha solo baja a $0.3155$: el
91\% es atribuible al cambio de ubicación. El colapso es específico de los
animales (\texttt{car} conserva AP@0.5 $= 0.9950$ en terreno nuevo mientras
\texttt{bird} cae de $0.5592$ a $0.0132$), y Grad-CAM sobre el backbone, con la
capa elegida por medición y no por convención, muestra que para animales que
ocupan menos del 1\% del cuadro la activación en ubicaciones nuevas está
\emph{menos} concentrada sobre el animal que un mapa uniforme. La robustez ante
oclusión sintética es nula: tapar el 25\% del área de cada animal cuesta el 69\%
del mAP. Sobre esas salidas se construye un módulo de monitoreo que agrupa las
detecciones en eventos independientes y reporta tasas por 100 noches-cámara; su
índice de nocturnidad reproduce la historia natural conocida de las catorce
especies detectadas sin haber recibido esa información, y al mismo tiempo, en
cámaras no vistas el conteo de eventos se desvía entre $-75\%$ y $+800\%$ según
la especie: sobreestima las comunes y pierde la más rara (\texttt{badger}).
